\clearpage
\chapter*{ABSTRACT}
\addcontentsline{toc}{chapter}{ABSTRACT}

\begin{center}
	\center
	\begin{singlespace}
		\large\bfseries\MakeUppercase{\thetitle}

		\normalfont\normalsize
		By:

		\bfseries \theauthor
	\end{singlespace}
\end{center}

\begin{singlespace}
	\small
	The transformation of electronic medical records (EMRs) requires access-control mechanisms that preserve patient control while supporting collaborative healthcare workflows. DecMed, as a decentralized EMR system, already uses IOTA, IPFS, and Proxy Re-Encryption (PRE); however, its previous mechanism still depended on patients to issue new access, remained coarse-grained because it relied mainly on role and purpose, and stored EMR data as large objects that were difficult to enforce according to care needs. This thesis aims to broaden DecMed access control so that access decisions cover actors, EMR objects, access rights, and access conditions more comprehensively, while also supporting data segmentation, limited delegation, revocation, auditing, and identity verification of actors on the delegation chain. The proposed approach applies Policy-Based Access Control (PBAC) as the policy model, Macaroon as a capability carrier with caveats for attenuation, off-chain EMR segmentation based on dataset categories and clinical functions, metadata and audit recording on IOTA, and PRE as the policy enforcement point before encrypted data is retrieved or reprocessed for an access recipient. The implementation includes EMR segmentation components, Macaroon-based tokens, authority enforcement on the PRE Server, IOTA Move smart contracts, the Patient Client, and the Hospital Client for workflows involving patients, medical record officers, nurses, physicians, laboratory officers, and pharmacists. Testing was conducted through functional scenarios \texttt{P-KF-01} to \texttt{P-KF-12} for outpatient and inpatient workflows, as well as nonfunctional tests \texttt{P-KNF-01} to \texttt{P-KNF-13}. The evaluation states that all formulated functional and nonfunctional requirements were fulfilled. The results show that policy-based access control, EMR segmentation, limited delegation, revocation, auditing, and identity verification of actors on the delegation chain can be operationalized in DecMed within the evaluated prototype implementation boundaries.

	\textbf{\textit{Keywords: electronic medical records, Policy-Based Access Control, Macaroon, IOTA, Proxy Re-Encryption}}
\end{singlespace}
\clearpage
